% HTN figure (top-down) -- include in a chapter with % HTN figure (top-down) -- include in a chapter with % HTN figure (top-down) -- include in a chapter with % HTN figure (top-down) -- include in a chapter with \input{chapters/htn_figure}
% Uses the compound/primitive/method styles and \htnpre from main.tex
\begin{figure}[htbp]
  \centering
  \begin{forest}
    for tree={
      font=\footnotesize, align=center, inner sep=3pt, rounded corners=2pt,
      minimum height=0.7cm,
      l sep=0.8cm, s sep=0.25cm,
      edge={draw=black!60, thick},
      parent anchor=south, child anchor=north,
    }
    [{\textbf{suppress-enemy-radar}}, compound
      [{m0: standard SEAD\htnpre{drone ready}}, method
        [{fly-to-\\area}, primitive]
        [{find-radar}, compound
          [{m1: passive\\search\htnpre{radar emitting}}, method
            [{listen-\\passively}, primitive]
          ]
          [{m2: decoy\\search\htnpre{radar silent}}, method
            [{act-as-\\decoy}, primitive]
            [{listen-\\passively}, primitive]
          ]
        ]
        [{suppress-\\radar}, primitive]
        [{check-\\result}, primitive]
        [{return-to-\\base}, primitive]
      ]
    ]
  \end{forest}

  \medskip
  % Legend
  \begin{tikzpicture}[font=\scriptsize,
      every node/.style={inner sep=2pt, rounded corners=2pt}]
    \node[draw, double, fill=blue!8] (c) {compound task};
    \node[draw, fill=white, right=6pt of c] (p) {primitive task};
    \node[draw, dashed, fill=orange!8, right=6pt of p] (m) {method (with precondition)};
  \end{tikzpicture}

  \caption{HTN for a SEAD mission. The compound task
    \textit{suppress-enemy-radar} is decomposed by method \textit{m0} into
    totally ordered subtasks (left to right). The compound task
    \textit{find-radar} has two alternative methods; the planner selects
    the one whose precondition holds in the current state.}
  \label{fig:htn-sead}
\end{figure}

% Uses the compound/primitive/method styles and \htnpre from main.tex
\begin{figure}[htbp]
  \centering
  \begin{forest}
    for tree={
      font=\footnotesize, align=center, inner sep=3pt, rounded corners=2pt,
      minimum height=0.7cm,
      l sep=0.8cm, s sep=0.25cm,
      edge={draw=black!60, thick},
      parent anchor=south, child anchor=north,
    }
    [{\textbf{suppress-enemy-radar}}, compound
      [{m0: standard SEAD\htnpre{drone ready}}, method
        [{fly-to-\\area}, primitive]
        [{find-radar}, compound
          [{m1: passive\\search\htnpre{radar emitting}}, method
            [{listen-\\passively}, primitive]
          ]
          [{m2: decoy\\search\htnpre{radar silent}}, method
            [{act-as-\\decoy}, primitive]
            [{listen-\\passively}, primitive]
          ]
        ]
        [{suppress-\\radar}, primitive]
        [{check-\\result}, primitive]
        [{return-to-\\base}, primitive]
      ]
    ]
  \end{forest}

  \medskip
  % Legend
  \begin{tikzpicture}[font=\scriptsize,
      every node/.style={inner sep=2pt, rounded corners=2pt}]
    \node[draw, double, fill=blue!8] (c) {compound task};
    \node[draw, fill=white, right=6pt of c] (p) {primitive task};
    \node[draw, dashed, fill=orange!8, right=6pt of p] (m) {method (with precondition)};
  \end{tikzpicture}

  \caption{HTN for a SEAD mission. The compound task
    \textit{suppress-enemy-radar} is decomposed by method \textit{m0} into
    totally ordered subtasks (left to right). The compound task
    \textit{find-radar} has two alternative methods; the planner selects
    the one whose precondition holds in the current state.}
  \label{fig:htn-sead}
\end{figure}

% Uses the compound/primitive/method styles and \htnpre from main.tex
\begin{figure}[htbp]
  \centering
  \begin{forest}
    for tree={
      font=\footnotesize, align=center, inner sep=3pt, rounded corners=2pt,
      minimum height=0.7cm,
      l sep=0.8cm, s sep=0.25cm,
      edge={draw=black!60, thick},
      parent anchor=south, child anchor=north,
    }
    [{\textbf{suppress-enemy-radar}}, compound
      [{m0: standard SEAD\htnpre{drone ready}}, method
        [{fly-to-\\area}, primitive]
        [{find-radar}, compound
          [{m1: passive\\search\htnpre{radar emitting}}, method
            [{listen-\\passively}, primitive]
          ]
          [{m2: decoy\\search\htnpre{radar silent}}, method
            [{act-as-\\decoy}, primitive]
            [{listen-\\passively}, primitive]
          ]
        ]
        [{suppress-\\radar}, primitive]
        [{check-\\result}, primitive]
        [{return-to-\\base}, primitive]
      ]
    ]
  \end{forest}

  \medskip
  % Legend
  \begin{tikzpicture}[font=\scriptsize,
      every node/.style={inner sep=2pt, rounded corners=2pt}]
    \node[draw, double, fill=blue!8] (c) {compound task};
    \node[draw, fill=white, right=6pt of c] (p) {primitive task};
    \node[draw, dashed, fill=orange!8, right=6pt of p] (m) {method (with precondition)};
  \end{tikzpicture}

  \caption{HTN for a SEAD mission. The compound task
    \textit{suppress-enemy-radar} is decomposed by method \textit{m0} into
    totally ordered subtasks (left to right). The compound task
    \textit{find-radar} has two alternative methods; the planner selects
    the one whose precondition holds in the current state.}
  \label{fig:htn-sead}
\end{figure}

% Uses the compound/primitive/method styles and \htnpre from main.tex
\begin{figure}[H]
  \centering
  \begin{forest}
    for tree={
      font=\footnotesize, align=center, inner sep=3pt, rounded corners=2pt,
      minimum height=0.7cm,
      l sep=0.8cm, s sep=0.25cm,
      edge={draw=black!60, thick},
      parent anchor=south, child anchor=north,
    }
    [{\textbf{suppress-enemy-radar}}, compound
      [{m0: standard SEAD\htnpre{drone ready}}, method
        [{fly-to-\\area}, primitive]
        [{find-radar}, compound
          [{m1: passive\\search\htnpre{radar emitting}}, method
            [{listen-\\passively}, primitive]
          ]
          [{m2: decoy\\search\htnpre{radar silent}}, method
            [{act-as-\\decoy}, primitive]
            [{listen-\\passively}, primitive]
          ]
        ]
        [{suppress-\\radar}, primitive]
        [{check-\\result}, primitive]
        [{return-to-\\base}, primitive]
      ]
    ]
  \end{forest}

  \medskip
  % Legend
  \begin{tikzpicture}[font=\scriptsize,
      every node/.style={inner sep=2pt, rounded corners=2pt}]
    \node[draw, double, fill=blue!8] (c) {compound task};
    \node[draw, fill=white, right=6pt of c] (p) {primitive task};
    \node[draw, dashed, fill=orange!8, right=6pt of p] (m) {method (with precondition)};
  \end{tikzpicture}

  \caption{HTN for a SEAD mission. The compound task
    \textit{suppress-enemy-radar} is decomposed by method \textit{m0} into
    totally ordered subtasks (left to right). The compound task
    \textit{find-radar} has two alternative methods; the planner selects
    the one whose precondition holds in the current state.}
  \label{fig:htn-sead}
\end{figure}
